\documentclass[
%a4paper,12pt
encoding=utf8
]{./twoeskd}

% \usepackage{eskdappsheet}

% Packages required by doxygen
\usepackage[export]{adjustbox} % also loads graphicx
\usepackage{graphicx}
\usepackage[utf8]{inputenc}
\usepackage{multicol}
\usepackage{multirow}
\usepackage{makeidx}

% NLS support packages
\usepackage[T2A]{fontenc}
\usepackage[russian]{babel}
\usepackage{pscyr}

% Font selection
\usepackage{courier}
\usepackage{amssymb}

\setlength{\parindent}{0cm}
\setlength{\parskip}{0.2cm}

% debug to see the frame borders
% from https://en.wikibooks.org/wiki/LaTeX/Page_Layout
% \usepackage{showframe}

% Indices & bibliography
\usepackage[backend=biber
           % ,style=authoryear-icomp
            ]{biblatex}
\addbibresource{used_sources.bib}
\usepackage[titles]{tocloft}
\setcounter{tocdepth}{3}
\setcounter{secnumdepth}{5}

% change style of titles in \section{}
\usepackage{titlesec}
\titleformat{\section}[hang]{\huge\bfseries\center}{\thetitle.}{1em}{}
\titleformat{\subsection}[hang]{\Large\normalfont\raggedright}{\thetitle.}{1em}{\underline}
\titleformat{\subsubsection}[hang]{\large\normalfont\raggedright}{\thetitle.}{1pt}{}

% Packages for text layout in normal mode
% \usepackage[parfill]{parskip} % автоматом делает пустые линии между параграфами, там где они есть в тексте
% \usepackage{indentfirst} % indent even in first paragraph
\usepackage{setspace}	 % controls space between lines
\setstretch{1} % space between lines
\setlength\parindent{0.9cm} % size of indent for every paragraph
\usepackage{csquotes}% превратить " " в красивые двойные кавычки
\MakeOuterQuote{"}


% this makes items spacing single-spaced in enumerations.
\newenvironment{my_enumerate}{
\begin{enumerate}
  \setlength{\itemsep}{1pt}
  \setlength{\parskip}{0pt}
  \setlength{\parsep}{0pt}}{\end{enumerate}
}


% Custom commands
% configure eskd
\titleTop{
	{\Large ПРАВИТЕЛЬСТВО РОССИЙСКОЙ ФЕДЕРАЦИИ \\
		НАЦИОНАЛЬНЫЙ ИССЛЕДОВАТЕЛЬСКИЙ УНИВЕРСИТЕТ \\
		«ВЫСШАЯ ШКОЛА ЭКОНОМИКИ» } \\
	\vspace*{0.2cm}
	{Факультет компьютерных наук \\
		Департамент программнoй инженерии \\
	}
}
\titleDesignedBy{Студент группы БПИ 151 НИУ-ВШЭ}{Куприянов К. И.}

\titleAgreedBy{
	\parbox[t]{7cm} {
		\centerline{Академический руководитель ОП}
		\centerline{Системная и программная инженерия}
		\centerline{профессор ДПИ}
}}{Д.В. Александров}
\titleApprovedBy{
	\parbox[t]{7cm} {
		\centerline{Академический руководитель}
		\centerline{образовательной программы}
		\centerline{«Программная Инженерия»}
		\centerline{профессор, канд. техн. наук}
}}{В. В. Шилов}
\titleName{Клиент-Серверное Android-Приложение для Управления Скидками в Розничных Сетях}
\workTypeId{RU.17701729.506900 12 01-1}

\titleSubname{Текст программы}


%===== C O N T E N T S =====


\begin{document}

% Titlepage & ToC
\pagenumbering{roman}

% some water filling text, that is pointless but adds text
% 
\tab[0.75cm] Техническое задание – это основной документ, оговаривающий набор требований и
порядок создания программного продукта, в соответствии с которым производится разработка
программы, ее тестирование и приемка.

Настоящее Техническое задание на разработку ``Клиент-Серверное Android-Приложение для Управления Скидками в Розничных Сетях'' содержит следующие разделы: ``Введение'', ``Основания для разработки'',
``Назначение разработки'', ``Требования к программе'', ``Требования к программным документам'',
``Технико-экономические показатели'', ``Стадии и этапы разработки'', ``Порядок контроля и
приемки'' и приложения.

В разделе ``Введение'' указано наименование и краткая характеристика области применения
``Клиент-Серверного Android-Приложения для Управления Скидками в Розничных Сетях''.

В разделе ``Основания для разработки'' указан документ на основании, которого ведется
разработка и наименование темы разработки.

В разделе ``Назначение разработки'' указано функциональное и эксплуатационное
назначение программного продукта.

Раздел ``Требования к программе'' содержит основные требования к функциональным
характеристикам, к надежности, к условиям эксплуатации, к составу и параметрам технических
средств, к информационной и программной совместимости, к маркировке и упаковке, к
транспортировке и хранению, а также специальные требования.

Раздел ``Требования к программным документам'' содержит предварительный состав
программной документации и специальные требования к ней.

Раздел ``Технико-экономические показатели'' содержит ориентировочную экономическую
эффективность, предполагаемую годовую потребность, экономические преимущества разработки
``Клиент-Серверного Android-Приложения для Управления Скидками в Розничных Сетях''.

Раздел ``Стадии и этапы разработки'' содержит стадии разработки, этапы и содержание
работ.

В разделе ``Порядок контроля и приемки'' указаны общие требования к приемке работы.

Настоящий документ разработан в соответствии с требованиями:\\
1) ГОСТ 19.101-77 Виды программ и программных документов \cite{gost_types_of_software};\\
2) ГОСТ 19.102-77 Стадии разработки \cite{gost_stages_of_devel};\\
3) ГОСТ 19.103-77 Обозначения программ и программных документов \cite{gost_marking_software};\\
4) ГОСТ 19.104-78 Основные надписи \cite{gost_main_signs};\\
5) ГОСТ 19.105-78 Требования к программным документам \cite{gost_demands_for_docs};\\
6) ГОСТ 19.201-78 Техническое задание. Требования к содержанию и оформлению \cite{gost_tz}.\\

Изменения к данному Техническому заданию оформляются согласно ГОСТ 19.603-78 \cite{gost_main_rules_change},
Перед прочтением данного документа рекомендуется ознакомиться с терминологией,
приведенной в Приложении 1 настоящего технического задания.


\newpage
\pagenumbering{arabic}
\tableofcontents

% --- add my custom headers ---
\newpage
\section{Текст программы}
Андроид приложение состоит из 19 классов, crawler - из 4.
Текст программы на исходном языке находится в архиве “Текст
программы.zip” в корневой директории архива в облачном хранилище по
адресу https://yadi.sk/d/JYE81i7x3UJJNP или по qr-коду (рис.~\ref{qr}) в связи с большим количеством строк
кода.

\begin{figure}[h!]
    \centering
    \includegraphics[width=0.7\textwidth]{./qr.png}
    \caption{ссылка на программу}
    \label{qr}
\end{figure}



\newpage
\section{Приложение 1. Список используемой литературы}
\subsection{Список используемой литературы}
\begin{my_enumerate}
    \item
    ГОСТ 19.102-77 Стадии разработки. //Единая система программной документации. -М.: ИПК Издательство стандартов, 2001.
    
    \item
    ГОСТ 19.201-78 Техническое задание. Требования к содержанию и оформлению // Единая система программной документации. -М.:ИПК Издательство стандартов, 2001.
    
    \item  ГОСТ 19.404-79 Пояснительная записка. Требования к содержанию и оформлению. //Единая система программной документации. – М.: ИПК Издательство стандартов, 2001
    
    \item
    ГОСТ 19.101-77 Виды программ и программных документов
    //Единая система программной документации. -М.: ИПК Издательство стандартов, 2.: 001.
    
    \item
    ГОСТ 19.103-77 Обозначения программ и программных документов. //Единая система программной документации. -М.: ИПК Издательство стандартов, 2001.
    
    \item
    ГОСТ 19.104-78 Основные надписи //Единая система программной документации. -М.: ИПК Издательство стандартов, 2001.
    
    \item 
    ГОСТ 19.106-78 Требования к программным документам, выполненным печатным способом. //Единая
    система программной документации. – М.: ИПК Издательство стандартов, 2001
    
    \item 
    ГОСТ 19.603-78 Общие правила внесения изменений. //Единая система программной документации. –
    М.: ИПК Издательство стандартов, 2001
    
    \item
    Oculus Documentation [Электронный ресурс]: Режим доступа: https:// developer3.oculus. com/documentation/
    
    \item
    Oculus Developers Blog [Электронный ресурс]: chrispruett – Squeezing Performance out of your Unity Gear VR Game, 2015 - Режим доступа: https://developer3.oculus.com/blog /squeezing-performance-out-of-your-unity-gear-vr-game/
    
    \item 
    Uninty Scripting Reference [Электронный ресурс]: Режим доступа: https: //docs.unity3d. com/ScriptReference/
    
\end{my_enumerate}



% Index
\newpage
\eskdListOfChanges

% \phantomsection
% \addcontentsline{toc}{section}{Алфавитный указатель}
% \printindex

\end{document}
